% Generated draft: requires mathematical and source review.
\documentclass[11pt]{article}
\usepackage[margin=25mm]{geometry}
\usepackage{fontspec}
\setmainfont{Latin Modern Roman}
\usepackage{amsmath,amssymb,mathtools,unicode-math}
\setmathfont{Latin Modern Math}
\usepackage{enumitem,needspace,xurl,hyperref}
\setlength{\parindent}{0pt}
\setlength{\parskip}{5pt}
\setlength{\emergencystretch}{3em}
\newcommand{\N}{\mathbb{N}}
\newcommand{\Z}{\mathbb{Z}}
\newcommand{\Q}{\mathbb{Q}}
\newcommand{\R}{\mathbb{R}}
\newcommand{\C}{\mathbb{C}}
\newcommand{\F}{\mathbb{F}}
\newcommand{\A}{\mathbb{A}}
\newcommand{\PP}{\mathbb{P}}
\newcommand{\E}{\mathbb{E}}
\newcommand{\eps}{\varepsilon}
\newcommand{\dd}{\,\mathrm{d}}
\begin{document}
\section*{UnsolvedMath ID: 20002325}
\textbf{Source problem code:} catalogue entry
\subsection{Definitions and mathematical statement}
Question 24 asks whether Borel completeness of \(\cong_T\) implies Borel completeness of \(\cong_{T'}\), but the excerpt omits the relation of \(T'\) to \(T\). The aid uses the one-named-constant interpretation: \(T\) is complete in a countable language, \(p\in S_1(T)\), and
\[
 T_p=\operatorname{Th}(M,a)\quad\text{for }M\models T,\ a\models p,
\]
with one new constant naming \(a\). This complete expansion is determined by \(p\); it is not the incomplete expansion allowing an arbitrary unnamed type for the constant.

Use standard spaces of countable models with domain \(\omega\). For equivalence relations \(E,F\), \(E\le_B F\) means there is a Borel map reducing equivalence in both directions. An isomorphism relation is Borel complete when every isomorphism relation on countable structures in a countable language Borel reduces to it, equivalently when countable graph isomorphism does. The term does not assert that the isomorphism relation itself is a Borel subset of the square.

Let \(\operatorname{Mod}(T)_p\) be the models realizing \(p\), and let \(E_p\) be pointed isomorphism for \(T_p\). Its countable-set Friedman--Stanley jump \(E_p^+\) identifies sequences exactly when they enumerate the same set of \(E_p\)-classes:
\[
 (x_i)\ E_p^+\ (y_j)
 \quad\Longleftrightarrow\quad
 (\forall i\,\exists j\ x_iE_py_j)\ \wedge\
 (\forall j\,\exists i\ x_iE_py_j).
\]
Multiplicity in the sequences is ignored. The aid reports
\[
 \cong_T|_{\operatorname{Mod}(T)_p}\ \le_B\ E_p^+.
\]
Its map lists all pointings \((M,a)\) for \(a\in p(M)\), with repetitions if only finitely many realize \(p\). The restriction to realizing models is essential; not every \(T\)-model need realize a nonisolated \(p\).

A stronger special case removes the jump. Suppose \(\mathcal C\subseteq\operatorname{Mod}(T)_p\) is an invariant Borel family on which \(\cong_T\) is Borel complete and \(\operatorname{Aut}(M)\) acts transitively on \(p(M)\) for every \(M\in\mathcal C\). Then Borel choice of one realization in each countable presentation gives
\[
 \cong_T|_{\mathcal C}\le_B E_p,
\]
so \(E_p\) is Borel complete. In particular, if a parameter-free formula defines a singleton in every \(T\)-model, naming that element gives Borel bireducible unpointed and pointed isomorphism relations. The aid does not claim that the transitivity/coding-family condition is automatic. Its general one-constant preservation question remains open in the supplied assessment, and Borel completeness of a jump is not used as a proof of Borel completeness of the underlying relation.
\subsection{Short English statement}
The aid interprets the question as naming one element of a fixed type. It reduces unpointed models realizing that type to sets of pointed isomorphism classes, and gets a direct reduction when the realizations form one automorphism orbit on a suitable coding family. Those conditions are special cases, not a resolution of unrestricted preservation.
\subsection{Sources}
\begin{itemize}
\item Local UnsolvedMath registry, permanent ID 20002325, exact source code catalogue entry; canonical \texttt{aim-logic-notes.json}, \texttt{notes[100]}, tag \texttt{question}; AIM Workshop Problem Lists, Logic, \emph{Computable stability theory}, Question 24. Created and updated 14 August 2026.
\item Source URL \url{https://aimath.org/pastworkshops/computestabproblems.pdf}. The exact one-line question gives no preceding definition of \(T'\). The PDF was not accessed.
\item The aid invokes the countable-set Friedman--Stanley jump without full bibliographic details or a separate theorem citation for the reduction.
\item Attempt 1: \texttt{reduction}, problem status at run \texttt{open}, rigor label \texttt{proved}, novelty confidence \texttt{low}; classification \texttt{PARTIAL-PROGRESS}, suggested difficulty L4, registry status \texttt{partially\_solved}. Candidate: realized-type jump reduction and orbit-transitive special case. General openness is a dated source assessment.
\item Ulam AI / UnsolvedMath Contributors; CC BY 4.0. Machine-generated research aid, not peer reviewed; independently verify before citation.
\end{itemize}
\end{document}
